\section{Unique fine cells outside the dummy region}
\label{sec:area_law_fine_cell_partition}

The half-open fine cells refine the dyadic layer partition of
\cite[Section~11]{OpenAI2026AreaLaw}. Retain the origin $o$, the finite
endpoint set $Z$, the integer width $C_0$, and the fine exponents $\ell_k$.

% Source: 10-geometry.tex, lines 154--181 and 200--207.
% These statements assign half-open fine cells. They do not assign the
% final primary and belt regions or prove the full two-family proposition.

\begin{lemma}[Unique fine cell outside the initial neighborhood]
    \label{thm:area_law_unique_fine_cell}
    \lean{TNLean.PEPS.AreaLaw.Geometry.exists_unique_fineCell_of_not_mem_dyadicNeighborhood}
    \leanok
    \uses{def:area_law_dyadic_layers,def:area_law_fine_layer_indices,
        def:area_law_fine_belt_scales}
    Suppose that $Z$ is nonempty and $C_0\ge2$. For every initial layer index
    $k_0$ and every $x\notin N_{k_0}$, there is a unique pair $(k,z)$
    with $k\ge k_0$, $z\in F_{k,\ell_k}$, and $x\in C_{\ell_k,z}$.
\end{lemma}
\begin{proof}\leanok
    \uses{thm:area_law_unique_layer,thm:area_law_fine_layer_count,
        thm:area_law_fine_cell_containment,thm:area_law_fine_belt_scales,
        thm:area_law_dyadic_indices}
    The point $x$ belongs to a unique layer $D_k$ with $k\ge k_0$. Since
    $\ell_k\le k$, the decomposition of $D_k$ into scale-$\ell_k$ fine cells
    gives $z\in F_{k,\ell_k}$ with $x\in C_{\ell_k,z}$. Conversely, if
    $h\ge k_0$, $w\in F_{h,\ell_h}$, and $x\in C_{\ell_h,w}$, then
    $C_{\ell_h,w}\subseteq D_h$, so $h=k$, and $w=c_{\ell_k}(x)=z$.
\end{proof}

\begin{lemma}[Exact cover by the dummy region and fine cells]
    \label{thm:area_law_dummy_fine_cells_cover}
    \lean{TNLean.PEPS.AreaLaw.Geometry.dyadicNeighborhood_union_iUnion_fineCells_eq_univ}
    \leanok
    \uses{def:area_law_dyadic_layers,def:area_law_fine_layer_indices,
        def:area_law_fine_belt_scales}
    For nonempty $Z$, integer $C_0\ge2$, and every initial layer index $k_0$,
    \begin{align}
        \mathbb R^2=N_{k_0}\cup\bigcup_{k\ge k_0}
        \bigcup_{z\in F_{k,\ell_k}}C_{\ell_k,z}.
    \end{align}
\end{lemma}
\begin{proof}\leanok
    \uses{thm:area_law_unique_fine_cell}
    Every $x\in\mathbb R^2$ either lies in $N_{k_0}$ or belongs to the fine
    cell $C_{\ell_k,z}$ of its unique pair $(k,z)$ with $k\ge k_0$ and
    $z\in F_{k,\ell_k}$.
\end{proof}
